%%
%% This is file `leaflet-manual.tex',
%% generated with the docstrip utility.
%%
%% The original source files were:
%%
%% leaflet.dtx  (with options: `manual')
%%
%% Copyright (C) 2003, 2004
%% Rolf Niepraschk, Rolf.Niepraschk@ptb.de
%% Hubert Gaesslein, HubertJG@open.mind.de
%%
%% This work may be distributed and/or modified under the
%% conditions of the LaTeX Project Public License, either version 1.3
%% of this license or (at your option) any later version.
%% The latest version of this license is in
%%   http://www.latex-project.org/lppl.txt
%% and version 1.3 or later is part of all distributions of LaTeX
%% version 2003/12/01 or later.
%%
%% This work has the LPPL maintenance status "author-maintained".
%%
\def\filename{leaflet-manual.tex}
\def\fileversion{v1.0d}   % change this when leaflet-manual changed, too.
\def\filedate{2012/06/04}
\def\docdate {2012/06/04} % change this when leaflet-manual changed, too.
\listfiles
\errorcontextlines=99
\documentclass[
%%notumble,
%%nofoldmark,
%%dvipdfm,
%%portrait,
%%titlepage,
%%nocombine,
%%a3paper,
%%debug,
%%nospecialtricks,
%%draft,
]{leaflet}


\renewcommand*\foldmarkrule{.3mm}
\renewcommand*\foldmarklength{5mm}

\usepackage[T1]{fontenc}
\usepackage{textcomp}
\usepackage{mathptmx}
\usepackage[scaled=0.9]{helvet}
\makeatletter
\def\ptmTeX{T\kern-.1667em\lower.5ex\hbox{E}\kern-.075emX\@}
\DeclareRobustCommand{\ptmLaTeX}{L\kern-.3em
        {\setbox0\hbox{T}%
         %\vb@xt@ % :-)
         \vbox to\ht0{\hbox{%
                            \csname S@\f@size\endcsname
                            \fontsize\sf@size\z@
                            \math@fontsfalse\selectfont
                            A}%
                      \vss}%
        }%
        \kern-.12em
        \ptmTeX}
\makeatother
\let\TeX=\ptmTeX
\let\LaTeX=\ptmLaTeX
\usepackage{shortvrb}
\MakeShortVerb{\|}
\usepackage{url}
\usepackage{graphicx}
\usepackage[dvipsnames,usenames]{color}
\usepackage{listings}
\definecolor{LIGHTGRAY}{gray}{.9}

%%%%\renewcommand{\descfont}{\normalfont}
\newcommand\Lpack[1]{\textsf{#1}}
\newcommand\Lclass[1]{\textsf{#1}}
\newcommand\Lopt[1]{\texttt{#1}}
\newcommand\Lprog[1]{\textit{#1}}

\newcommand*\defaultmarker{\textsuperscript\textasteriskcentered}

\title{Processing Kickstart\vspace{-2ex}}
\author{%
  Dan Lidral-Porter\vspace{-2ex}
}
\date{}

\CutLine*{1}
\CutLine*{6}

%AddToBackground{5}{%  Background of a small page
% \put(0,0){\textcolor{Cerulean}{\rule{\paperwidth}{\paperheight}}}}
%
%AddToBackground*{2}{% Background of a large page
% \put(\LenToUnit{.5\paperwidth},\LenToUnit{.5\paperheight}){%
%   \makebox(0,0)[c]{%
%     \resizebox{.9\paperwidth}{!}{\rotatebox{35.26}{%
%       \textsf{\textbf{\textcolor{LIGHTGRAY}{BACKGROUND}}}}}}}}

\begin{document}

\definecolor{framegrey}{rgb}{0.7,0.7,0.7}

\lstset{language=Java
       ,basicstyle=\ttfamily
       ,frame=single
       ,rulecolor=\color{framegrey}
       }

\maketitle
\thispagestyle{empty}

%%\LARGE

\section{The Canvas}

Processing artworks are called \textit{sketches}.
Every sketch begins with defining the size of the window it will use, called the \textit{canvas}.

The \texttt{size} command is used to define the size of the canvas, in pixels.
The general form of the size command is \texttt{size(width, height)}.
The values for width and height must be whole numbers (no decimal points).

Here's an example:

\begin{lstlisting}
size(1280, 720);
\end{lstlisting}

This creates a canvas that is 1280 pixels wide and 720 pixels high.
Try creating a sketch with only a size command, and change the numbers to experiment with the width and height.

A location within the canvas is described using two numbers, called the \textit{coordinates}.
This figure illustrates the names of the coordinates:

\begin{figure}[!h]
  \centering
  \includegraphics[scale=.25]{canvas.eps}
\end{figure}

Coordinates are displayed as \texttt{(x, y)}, where \texttt{x} is the x-coordinate, and \texttt{y} is the y-coordinate.

\textbf{Exercise:} where would the red dot be if the x-coordinate were zero?
What about the y-coordinate?
What about both?

\section{Color}

Computers store colors as amount of red, green, and blue light, but this is hard to reason about.
A more natural way of representing colors is with \textit{hue}, \textit{saturation}, and \textit{brightness}.

The \textit{hue} is the base coloring of the color, before it's modified by the saturation and brightness values.
The hue starts and ends at red, and in between covers the whole color wheel.

The \textit{saturation} is how strong the color is.
Full saturation means the color is as intensely colored as possible.
A saturation of zero produces values of grey.

The \textit{brightness} is how much light the color emits.
A brightness of zero means the color is black, no matter what the other values.

\textbf{Exercise:} consider pure white. What is its brightness? How about its saturation? Does its hue matter?

You can tell Processing to use hue, saturation, and brightness by using the \texttt{colorMode} command:

\begin{lstlisting}
colorMode(HSB, 100);
\end{lstlisting}

The \texttt{HSB} means you want to use Hue, Saturation, and Brightness mode, and the \texttt{100} means that you want the maximum value for each to be 100.
You should put this in a sketch right after the \texttt{size} command.

After the \texttt{colorMode} command, you can create colors by using the \texttt{color} command.
The general form of the color command is \texttt{color(hue, saturation, brightness)}.

A simple thing to use colors with is the \texttt{background} command, which paints the entire canvas with the color you give it.
Here is a sketch that paints the whole canvas yellow using the \texttt{background} command:

\begin{lstlisting}
size(1280, 720);
colorMode(HSB, 100);
background( color(14, 100, 90) );
\end{lstlisting}

\textbf{Exercise:} write a sketch that paints the canvas purple.

\section{Shapes}

\section{All Together}


\begin{thebibliography}{000}
\bibitem{cit:processing-ref}
  \textsc{B. Fry, C. Reas, et. al.}: The Processing Language Reference. 2001-\\
  \url{http://processing.org/reference}
\end{thebibliography}

\loggingall
\end{document}
\endinput
%%
%% End of file `leaflet-manual.tex'.
